% ============================================================
% f2014_shaanxi_arts_q17_tetrahedron_views.tex
% 2014 陕西卷（文科）第 17 题：四面体 ABCD 及其三视图
% 来源 PDF：普通高考/2014/2014陕西文.pdf
% 原页：第 1 页（A3），印面 913，第 17 题题干下方
% 由三视图得 DA⊥DB、DA⊥DC、DB⊥DC，且 DA=1、DB=DC=2；
% 原页俯视图只标出纵向尺寸 2，横向尺寸由主视图的 2 得出。
% 左图为立体示意图（平行投影，不按比例）；右图为三视图。
% ============================================================
\begin{tikzpicture}[>=Stealth, line width=0.7pt, font=\small,
  line cap=round, line join=round]
  % ---------- 立体示意图 ----------
  \begin{scope}[scale=1.28]
  \coordinate (A) at (0.10,2.40);
  \coordinate (B) at (-1.50,0.35);
  \coordinate (C) at (2.60,0.75);
  \coordinate (D) at (0.05,0.95);
  \coordinate (E) at (-0.70,1.375); % AB 中点
  \coordinate (F) at (-0.725,0.65); % BD 中点
  \coordinate (G) at (1.325,0.85);  % DC 中点
  \coordinate (H) at (1.35,1.575);  % CA 中点

  % 可见棱
  \draw (A) -- (B) -- (C) -- cycle;
  \draw (E) -- (H);
  % 不可见棱与截面边
  \draw[dashed] (A) -- (D) -- (B);
  \draw[dashed] (D) -- (C);
  \draw[dashed] (E) -- (F) -- (G) -- (H);

  \fill (A) circle (1.1pt) node[above] {$A$};
  \fill (B) circle (1.1pt) node[below left] {$B$};
  \fill (C) circle (1.1pt) node[right] {$C$};
  \fill (D) circle (1.1pt) node[below] {$D$};
  \fill (E) circle (1.1pt) node[left] {$E$};
  \fill (F) circle (1.1pt) node[below left] {$F$};
  \fill (G) circle (1.1pt) node[below] {$G$};
  \fill (H) circle (1.1pt) node[above right] {$H$};
  \end{scope}

  % ---------- 三视图 ----------
  % 主视图：横向 2，纵向 1
  \begin{scope}[shift={(5.05,0.25)}]
    \draw (0,1) -- (0,0) -- (2,0) -- cycle;
    \draw[<->] (-0.42,0) -- (-0.42,1);
    \node[left] at (-0.42,0.5) {$1$};
    \draw[<->] (0,-0.42) -- (2,-0.42);
    \node[below] at (1,-0.42) {$2$};
    \node[below] at (1,-0.85) {主视图};
  \end{scope}

  % 左视图：原页不标尺寸
  \begin{scope}[shift={(7.85,0.25)}]
    \draw (0,1) -- (0,0) -- (2,0) -- cycle;
    \node[below] at (1,-0.45) {左视图};
  \end{scope}

  % 俯视图：纵向 2，横向由主视图给出
  \begin{scope}[shift={(5.05,-4.20)}]
    \draw (0,0) -- (0,2) -- (2,2) -- cycle;
    \draw[<->] (-0.42,0) -- (-0.42,2);
    \node[left] at (-0.42,1) {$2$};
    \node[below] at (1,-0.45) {俯视图};
  \end{scope}
\end{tikzpicture}
